\section{Kết luận và kiến nghị}

\subsection{Kết luận}

Đề án cá nhân đã hoàn thành việc nghiên cứu, luận giải và giải quyết một cách toàn diện bốn chuẩn đầu ra của học phần Quản lý chất lượng giáo dục nghề nghiệp (mã học phần ED6111) thông qua trường hợp nghiên cứu điển hình tại Trường Cao đẳng Kỹ Nghệ II:

Về Chuẩn đầu ra 1, đề án đã phân tích thấu đáo bối cảnh tự chủ tài chính toàn diện, vận dụng thành công các công cụ quản trị chiến lược PESTEL và SWOT để làm rõ các lực đẩy môi trường; phân tích sâu sắc kết quả tự đánh giá đạt điểm tuyệt đối năm 2025 và chỉ ra ba khoảng trống chất lượng bản chất về tư duy hao hụt người học, độ trễ khảo sát và phân mảnh dữ liệu; từ đó xác lập bốn định hướng cải tiến chiến lược mang tính đột phá cho giai đoạn 2026 -- 2030.

Về Chuẩn đầu ra 2, đề án đã xây dựng khung ma trận tích hợp toàn diện giữa tám tiêu chí, năm mươi tiêu chuẩn của Thông tư số 14/2024/TT-BLĐTBXH với sáu mươi lăm quy trình vận hành nội bộ của trường; thiết kế hoàn chỉnh hai công cụ đo lường chuyên biệt gồm Phiếu khảo sát phản hồi người học theo thang đo Likert năm mức độ và Rubric đánh giá kỹ năng nghề tích hợp năng lực số; đồng thời thực hành biên soạn mẫu mực Báo cáo tự đánh giá cho Tiêu chuẩn 8.1 với đầy đủ minh chứng và kế hoạch hành động khắc phục sâu sắc.

Về Chuẩn đầu ra 3, đề án đã kiến tạo kiến trúc hệ thống đảm bảo chất lượng bên trong theo mô hình Quản lý chất lượng toàn diện thúc đẩy bởi công nghệ số tích hợp tiêu chuẩn quốc tế ISO 21001:2018; phân định ma trận trách nhiệm RACI rành mạch; đóng kín chu trình cải tiến Hoạch định -- Thực hiện -- Kiểm tra -- Cải tiến; thiết kế chi tiết hai quy trình thao tác chuẩn nòng cốt là Hệ thống cảnh báo sớm biến động sĩ số và Khảo sát phản hồi tức thời; đồng thời xây dựng chính sách hòa nhập cho người học có nhu cầu đặc biệt.

Về Chuẩn đầu ra 4, đề án đã phân tích và giải quyết thuyết phục một tình huống xung đột đạo đức nghề nghiệp điển hình về áp lực làm đẹp số liệu việc làm sinh viên tốt nghiệp; đề xuất Quy trình năm bước giải quyết xung đột dựa trên bốn trụ cột đạo đức: Khách quan, Trung thực, Liêm chính học thuật và Trách nhiệm xã hội; từ đó thiết lập hai cơ chế thể chế bảo đảm tính trung thực dữ liệu gồm Quy chế bảo vệ người phản ánh sai phạm và Cơ chế kiểm toán dữ liệu độc lập.

\subsection{Kiến nghị}

Nhằm hiện thực hóa các định hướng cải tiến chất lượng và lan tỏa mô hình quản trị hiệu quả, đề án đề xuất một số kiến nghị cụ thể:

Đối với Bộ Lao động -- Thương binh và Xã hội cùng Tổng cục Giáo dục nghề nghiệp, kiến nghị cơ quan quản lý sớm hoàn thiện và đưa vào vận hành Cổng cơ sở dữ liệu kiểm định chất lượng giáo dục nghề nghiệp quốc gia liên thông, cho phép tích hợp trực tiếp dữ liệu số hóa từ các cơ sở đào tạo nhằm tự động hóa công tác đối soát minh chứng; đồng thời ban hành cẩm nang hướng dẫn chuyên môn về bảo đảm liêm chính học thuật và đạo đức nghề nghiệp trong công tác tự đánh giá và đánh giá ngoài.

Đối với Hội đồng trường và Ban Giám hiệu Trường Cao đẳng Kỹ Nghệ II, kiến nghị nhà trường sớm xem xét phê duyệt chỉ tiêu duy trì sĩ số toàn khóa tối thiểu 88\% trong Bảng Mục tiêu chất lượng năm 2026; cấp kinh phí triển khai phân hệ Hệ thống cảnh báo sớm và công cụ khảo sát tức thời trên ứng dụng di động nội bộ của sinh viên ngay trong năm học 2026; đồng thời ban hành Quy chế bảo vệ người phản ánh sai phạm và thành lập Tổ Kiểm toán dữ liệu chất lượng độc lập trực thuộc Hội đồng trường.
